\chapter{Discussion}

\section{Interpretation of Results}

The final results indicate that the neural network can approximate the Black-Scholes pricing function with high accuracy in the sampled parameter domain.
This supports the main research question in the controlled synthetic setting considered in the project.

The comparison with Monte Carlo should be interpreted carefully.
The neural network is evaluated on the full held-out test set, whereas Monte Carlo is evaluated on a deterministic subset due to computational cost.
The purpose of the comparison is therefore not to declare a universally superior method, but to clarify the behavior of a learned surrogate relative to a classical simulation-based estimator.

\section{When the Neural Network Works Well}

The neural network works well when test samples come from the same parameter ranges used during training.
In this regime, the mapping from option parameters to Black-Scholes prices is smooth and deterministic.
The engineered moneyness feature and the SiLU activation function further improve the approximation in the experiments.

The runtime benchmark also highlights a practical advantage of the surrogate approach.
After training, neural inference is substantially faster than Monte Carlo simulation for repeated pricing queries in the benchmark setting.

\section{Sources of Error}

Prediction errors can arise from finite training data, imperfect optimization, model capacity limits, and numerical scaling choices.
Errors may also vary across the parameter space.
For this reason, the project includes diagnostic plots against moneyness, maturity, and volatility.

The noisy-target experiments show another source of error: imperfect labels.
When training targets are perturbed, the model remains accurate under moderate noise, but approximation quality decreases as the noise level increases.
This behavior is expected because the model receives less precise supervision.

\section{Limits of the Synthetic Setting}

The use of synthetic data is appropriate for the research question, because the project aims to approximate the Black-Scholes pricing function itself.
However, this also limits the scope of the conclusions.

Real exchange-traded option prices reflect additional effects such as bid-ask spreads, liquidity, dividends, supply and demand, and implied-volatility structure.
Training on such data would change the nature of the problem from approximating Black-Scholes to modeling market prices.
Therefore, the results should not be interpreted as evidence that the model can price real market options.

\section{Comparison with Classical Methods}

The analytical Black-Scholes formula remains the exact and fastest method in this specific setting.
Monte Carlo is useful as a numerical benchmark and as a method that generalizes to settings without closed-form solutions.
SVR provides a classical machine learning baseline on reduced datasets.

The neural network occupies a different role.
It is not intended to outperform the closed-form formula when the formula is available.
Its value lies in demonstrating how a supervised model can learn a known pricing map and act as a fast surrogate after training.
